\section{Results}
\label{sec:results}

The simulated methods produced similar central values
(\cref{tab:summary}). Method B was higher on average by
\SI{0.18}{\milli\metre}. The display in \cref{fig:comparison} makes both paired
proximity and the simulated status visible.

\begin{table}[htbp]
  \centering
  \caption{Summary of the simulated paired measurements. Values are mean
  (standard deviation).}
  \label{tab:summary}
  \begin{tabular}{lS[table-format=2.2]S[table-format=1.2]}
    \toprule
    Quantity & {Estimate} & {SD} \\
    \midrule
    Method A (\si{\milli\metre}) & 12.44 & 1.31 \\
    Method B (\si{\milli\metre}) & 12.62 & 1.35 \\
    Paired difference (\si{\milli\metre}) & 0.18 & 0.24 \\
    \bottomrule
  \end{tabular}
\end{table}

\begin{figure}[htbp]
  \centering
  \includegraphics[width=.78\linewidth]{method_comparison.pdf}
  \caption{Comparison of two simulated measurement methods. The diagonal is
  the line of equality; it is a visual guide rather than a fitted model.}
  \label{fig:comparison}
\end{figure}
